\section{Zielsetzung und Erkenntnisinteresse}
\label{sec:zielsetzung-und-erkenntnisinteresse}

Das Ziel dieser Bachelorarbeit ist es, das Two-Step-Flow-Modell und das Konzept von Opinion-Leadership zu analysieren, um die Frage zu beantworten, wie diese auf den heutigen Informationsfluss in sozialen Medien angewendet werden können. Entworfen wurde das Two-Step-Flow-Modell 1948 von Lazarsfeld, Berelson sowie Gaudet und wird hier auf den Mikroblogging-Dienst Twitter angewandt.

Das Tweetverhalten von US-Präsident Donald Trump wird untersucht und anhand einer Case-Study mit der klassischen Kommunikationstheorie verglichen, um zu ermitteln, ob die Theorie auf soziale Netzwerke angewendet werden kann.

Folgende Punkte sind die Ziele der Arbeit:

\begin{itemize}
    \item Die Studienergebnisse zu Trumps Tweetverhalten zusammenfassen,
    \item den Fall Nordstrom beschreiben und analysieren (Nachzeichnung des Two-Step-Flows),
    \item und anhand der Ergebnisse die folgende Forschungsfrage beantworten: Inwiefern erfüllt Donald Trumps Twittergebrauch die Kriterien von virtueller Meinungsführerschaft?
\end{itemize}

Es wird erwartet, dass gemäß den Forschungsergebnissen die Meinung in der Bevölkerung gegenüber Trump verglichen mit anderen Opinion-Leadern durchschnittlich negativer behaftet ist. Außerdem ist damit zu rechnen, dass seine Opinion-Followers aus einer Personengruppe stammen, die Trump wählen.